\section{Oversight: am I still in control?}
\label{sec:oversight}

The clinician must remain the responsible agent. The fourth question asks whether
the human stays in control: how to override, pause, or stop the system, and when it
escalates a decision to a qualified person rather than acting alone. The goal is
calibrated trust, meaning reliance proportioned to the system's demonstrated
capability, which sits between two failure modes that the oversight design must
hold apart.
Over-trust, or automation bias, is the tendency to defer to a machine even when it
is wrong; the clinical-decision-support literature documents how readily it occurs
when a confident recommendation appears on screen \cite{goddard2012automation}.
Under-trust, or algorithm aversion, is the opposite error of discarding useful
guidance. Appropriate reliance is the calibrated middle, and it depends on the
human keeping levers that are real rather than ceremonial \cite{lee2004trust}.

The framework provides those levers through two complementary supervision modes.
Human-in-the-loop means the clinician sits inside the decision: the system proposes
an action under verification before generation, but it cannot act until the
clinician accepts it, so a human authorizes each step. Human-on-the-loop means the
clinician supervises a running process from above: the system acts under audit
within pre-authorized bounds while the clinician monitors and retains the power to
intervene at any moment. The ten-gate VVUQ examination decides which mode applies
to a given action and, when a proposal falls outside its bounds, takes the ESCALATE
path that routes the decision to a qualified human.

Above both modes sits an unconditional control: the clinician can override and stop
the system at any time, for any reason, without satisfying the gates first.
\autoref{tab:oversight} sets out the four oversight positions, where the clinician
sits in each, and what control remains available. Figure 6 shows the same loop, in
which a monitored recommendation reaches a decision point that can resolve to
execute, escalate, or override. The design intent is that the human is never a
spectator: in-the-loop the human authorizes, on-the-loop the human supervises, on
escalation the human is summoned, and on override the human stops action.

\clearpage

This is how the framework guards against both automation bias and algorithm
aversion at once. The gates and the audit record give the clinician enough evidence
to rely on the system when it has earned reliance, which counters aversion; the
ever-present override and the escalation path ensure the clinician can refuse or
halt it the moment it has not, which counters over-trust. Control, defined this way,
is what makes calibrated trust possible rather than merely aspirational.

\begin{cfwfig}
\node[cftwo] (a) at (0,0) {Monitor}; 
\node[cfact] (b) at (4,0) {Recommend}; 
\node[cfdec] (c) at (7.3,0) {Decision};
\node[cfhope] (d) at (11.25,1.2) {Execute}; 
\node[cfthree] (e) at (11.25,0) {Escalate}; 
\node[cfrisk] (f) at (11.25,-1.2) {Override};
\draw[cfedge] (a)--(b); 
\draw[cfedge] (b)--(c);
\draw[cfedge] (c.east)--(d.west); 
\draw[cfedge] (c.east)--(e.west); 
\draw[cfedge] (c.east)--(f.west);
\end{cfwfig}
\vspace{-0.75cm}
\cfwcaption{Figure 6. The oversight loop: a
monitored recommendation reaches a decision that resolves to execute under audit,
escalate to a qualified human, or override and stop.}

\needspace{9\baselineskip}\begin{table}[H]
\footnotesize
\begin{tabularx}{\textwidth}{@{}L{2.7cm} Y L{5.2cm}@{}}
\toprule
\textbf{Oversight mode} & \textbf{Where the clinician sits} & \textbf{Control available}\\
\midrule
Human-in-the-loop & Inside the decision; the action waits for the clinician & Authorize, modify, or refuse each step\\
Human-on-the-loop & Above a running process within pre-authorized bounds & Monitor in real time, intervene at will\\
Escalation & Summoned when a proposal exceeds the gates' bounds & Decide the escalated case as a human\\
Override and stop & Outside the gates, at any moment & Halt system unconditionally, no gates\\
\bottomrule
\end{tabularx}
\caption{The four oversight positions: where the clinician sits and what control
each preserves, from human-in-the-loop authorization to an unconditional stop.}
\label{tab:oversight}
\end{table}

Control is a set of switches: a pause, an override, and an
escalation path that the clinician can reach at any moment. Those switches, written
into the illustrative H. R. 9510 as standing provisions rather than vendor
courtesies \cite{Kawchak2026HR9510Bill}, are the concrete answer to the question of
whether the clinician is still in control, and the mechanism that keeps trust
calibrated between automation bias and algorithm aversion.
